\rtmtight
\begin{tblr}{width=\textwidth,colspec={|Q[c,m,2.1cm]|X[l,m]|},hlines,vlines,hspan=minimal,rowsep=1.5pt,colsep=3pt,row{1}={bg=headgray,font=\bfseries}}
\SetCell{c,m}\hfil\logo\hfil & \textbf{POLITEKNIK NEGERI MALANG}\par\textbf{JURUSAN TEKNOLOGI INFORMASI}\par\textbf{PROGRAM STUDI: D4 TEKNIK INFORMATIKA} \\
\SetCell[c=2]{l} \textbf{RENCANA TUGAS MAHASISWA (RTM) DAN RUBRIK PENILAIAN} & \\
Nama Mata Kuliah & Pancasila \\
Kode Mata Kuliah & RTI251001 \quad \textbf{sks / Jam perminggu:} 2 SKS/2 jam \quad \textbf{Semester:} 1 \\
Dosen Pengampu & \begin{enumerate}\item Dr. Widaningsih Condrowardani, S.Psi., S.H., M.H.\end{enumerate} \\
\end{tblr}
\assessmentblock{Tugas Kajian Nilai Pancasila}{Tugas individu/kelompok}{SCPMK0102-00101, SCPMK0102-00102, SCPMK0102-00103, SCPMK0102-00104}{Mahasiswa menyusun kajian tertulis dan presentasi tentang materi Pancasila tiap pertemuan: tinjauan historis-filosofis, UUD 1945 dan kelembagaan negara, ideologi nasional, serta Pancasila dan HAM.}{Membaca modul dan pustaka, merangkum konsep, mengaitkan konsep dengan contoh aktual di masyarakat, menuliskan hasil kajian, dan mempresentasikan atau menjawab tanya jawab lisan.}{Ringkasan/makalah singkat, slide presentasi, dan/atau jawaban tes lisan terdokumentasi.}{Ketepatan konsep dan landasan Pancasila & 8 & Konsep, definisi, dan landasan historis-kultural-yuridis-filosofis dijelaskan benar, lengkap, dan merujuk sumber yang sahih. & Sebagian besar konsep benar, tetapi ada definisi atau landasan yang kurang lengkap atau tanpa rujukan. & Konsep banyak yang keliru, tertukar, atau tidak merujuk sumber sama sekali. \\ \\
Kualitas analisis keterkaitan nilai dengan kehidupan & 7 & Kajian mengaitkan nilai Pancasila dengan contoh nyata kehidupan bermasyarakat/bernegara secara logis dan relevan. & Keterkaitan nilai dengan contoh nyata ada, tetapi sebagian contoh kurang relevan atau dangkal. & Kajian hanya menyalin materi tanpa keterkaitan dengan kehidupan nyata. \\ \\
Kerapian dokumentasi dan penyampaian & 5 & Tulisan terstruktur, bahasa baku, presentasi lancar, dan pertanyaan dijawab dengan jelas dan sesuai materi. & Tulisan atau presentasi cukup jelas, tetapi struktur, bahasa, atau jawaban tanya jawab masih kurang konsisten. & Tulisan tidak terstruktur, presentasi tidak siap, atau pertanyaan tidak terjawab. \\}

\assessmentblock{Kuis Konsep Pancasila dan Ideologi}{Kuis}{SCPMK0102-00101, SCPMK0102-00103}{Mahasiswa mengerjakan Kuis 1 (konsep dasar dan filsafat Pancasila, minggu 3) dan Kuis 2 (ideologi nasional dan ideologi dunia, minggu 10).}{Mempelajari materi pertemuan sebelumnya, mengerjakan soal kuis tertulis secara mandiri dalam waktu terbatas, dan menuliskan jawaban ringkas disertai alasan.}{Lembar jawaban kuis tertulis (esai singkat dan/atau pilihan ganda).}{Penguasaan konsep dasar, filsafat, dan ideologi Pancasila & 5 & Jawaban menunjukkan penguasaan definisi, unsur filsafat, hierarki nilai, dan fungsi ideologi Pancasila dengan benar. & Sebagian jawaban benar, tetapi ada konsep inti yang tertukar atau tidak dijelaskan. & Sebagian besar jawaban keliru atau hanya menebak tanpa dasar konsep. \\ \\
Ketepatan penerapan konsep pada contoh & 3 & Konsep diterapkan tepat pada contoh kasus kehidupan sehari-hari atau perbandingan ideologi yang diminta soal. & Penerapan konsep pada contoh cukup tepat, tetapi belum lengkap atau kurang spesifik. & Jawaban tidak menerapkan konsep pada contoh atau penerapannya salah. \\ \\
Ketuntasan dan kejelasan jawaban & 2 & Seluruh soal dijawab tuntas dengan uraian ringkas, runtut, dan mudah dipahami. & Sebagian soal dijawab tuntas, tetapi ada uraian yang tidak runtut atau menggantung. & Banyak soal tidak dijawab atau uraian tidak dapat dipahami. \\}

\assessmentblock{Case Method Isu HAM, Korupsi, dan Pembangunan}{Case Method/FGD}{SCPMK0102-00104}{Mahasiswa menganalisis kasus aktual pelaksanaan HAM, tindak pidana korupsi, dan paradigma pembangunan berbasis nilai Pancasila secara berkelompok pada minggu 12--16.}{Membaca kasus, mengidentifikasi masalah dan dasar hukum/nilai yang relevan, berdiskusi kelompok (FGD), merumuskan sikap dan solusi, lalu mempresentasikan hasil analisis di kelas.}{Lembar analisis kasus, notulen diskusi kelompok, dan bahan presentasi hasil analisis.}{Ketajaman analisis kasus berbasis nilai Pancasila & 6 & Masalah kasus diidentifikasi tepat, dikaitkan dengan pasal/dasar hukum dan nilai sila yang relevan, serta didukung fakta kasus. & Masalah utama teridentifikasi, tetapi kaitan dengan dasar hukum atau nilai sila masih kurang spesifik. & Analisis tidak menyentuh masalah utama kasus atau tidak mengaitkan dengan nilai Pancasila. \\ \\
Kualitas argumentasi dan partisipasi diskusi & 5 & Argumen logis, santun, didukung data, dan mahasiswa aktif merespons pendapat anggota lain dalam FGD. & Argumen cukup logis, tetapi partisipasi pasif atau sebagian pendapat tanpa dukungan data. & Tidak berargumen, hanya mengulang pendapat orang lain, atau tidak berpartisipasi dalam diskusi. \\ \\
Rekomendasi solusi dan refleksi sikap & 4 & Rekomendasi pencegahan/solusi konkret, dapat dilaksanakan, dan disertai refleksi sikap pribadi sesuai nilai Pancasila. & Rekomendasi relevan tetapi masih umum, atau refleksi sikap belum menunjukkan internalisasi nilai. & Tidak ada rekomendasi yang operasional atau refleksi sikap tidak sesuai konteks kasus. \\}

\assessmentblock{UTS Pancasila dan Konstitusi}{Ujian tertulis (esai dan/atau pilihan ganda)}{SCPMK0102-00101, SCPMK0102-00102}{Mahasiswa mengerjakan ujian tengah semester tertulis mencakup materi minggu 1--7: tinjauan historis-kultural-yuridis-filosofis, sejarah perjuangan bangsa, sistem filsafat, UUD RI 1945, amandemen, Trias Politika, dan kelembagaan negara.}{Mempelajari seluruh materi minggu 1--7, mengerjakan soal esai dan/atau pilihan ganda secara mandiri dalam waktu ujian yang ditentukan.}{Lembar jawaban ujian tertulis.}{Penguasaan tinjauan historis, kultural, yuridis, dan filosofis Pancasila & 8 & Jawaban menjelaskan definisi, sejarah perjuangan bangsa, dan sistem filsafat Pancasila secara benar dan lengkap sesuai soal. & Jawaban benar pada sebagian besar soal, tetapi ada uraian sejarah atau filsafat yang tidak lengkap. & Jawaban banyak yang salah atau tidak menjawab inti pertanyaan. \\ \\
Pemahaman UUD 1945, amandemen, dan kelembagaan negara & 8 & Jawaban menguraikan UUD 1945, proses amandemen, Trias Politika, dan lembaga negara utama dengan tepat dan sesuai dasar hukum. & Sebagian uraian tepat, tetapi ada kekeliruan pada proses amandemen atau fungsi lembaga negara. & Uraian tentang konstitusi dan kelembagaan negara sebagian besar keliru. \\ \\
Kejelasan argumentasi jawaban esai & 4 & Jawaban esai runtut, menggunakan istilah yang tepat, dan disertai alasan atau contoh pendukung. & Jawaban esai cukup jelas, tetapi alasan atau contoh pendukung terbatas. & Jawaban esai tidak runtut, tanpa alasan, atau tidak dapat dipahami. \\}

\assessmentblock{UAS Ideologi, HAM, dan Pembangunan}{Ujian tertulis (esai dan/atau pilihan ganda)}{SCPMK0102-00103, SCPMK0102-00104}{Mahasiswa mengerjakan ujian akhir semester tertulis mencakup materi minggu 9--16: ideologi nasional, ideologi lain di dunia, Pancasila dan HAM, pelaksanaan HAM dalam UUD RI 1945, tindak pidana korupsi, dan Pancasila sebagai paradigma pembangunan.}{Mempelajari seluruh materi minggu 9--16, mengerjakan soal esai dan/atau pilihan ganda secara mandiri dalam waktu ujian yang ditentukan.}{Lembar jawaban ujian tertulis.}{Penguasaan ideologi nasional dan perbandingan ideologi dunia & 13 & Jawaban menjelaskan fungsi dan proses terbentuknya ideologi Pancasila serta membandingkannya dengan ideologi lain secara benar. & Penjelasan ideologi sebagian benar, tetapi perbandingan dengan ideologi lain kurang tepat atau dangkal. & Jawaban tidak membedakan Pancasila dengan ideologi lain atau sebagian besar keliru. \\ \\
Analisis HAM, antikorupsi, dan paradigma pembangunan & 14 & Jawaban menganalisis pelaksanaan HAM dalam UUD 1945, tindak pidana korupsi, dan paradigma pembangunan dengan dasar hukum dan contoh yang tepat. & Analisis cukup tepat, tetapi dasar hukum atau contoh pada sebagian jawaban tidak lengkap. & Jawaban tidak menganalisis, hanya menyebut istilah, atau dasar hukumnya salah. \\ \\
Logika dan koherensi argumentasi jawaban & 8 & Seluruh jawaban benar dan logis, uraian runtut, dan kesimpulan konsisten dengan argumen yang dibangun. & Jawaban umumnya logis, tetapi ada bagian yang tidak konsisten atau lompatan argumen. & Jawaban tidak logis, saling bertentangan, atau tidak menjawab pertanyaan. \\}

\begin{tblr}{width=\textwidth,colspec={|Q[l,m,3.2cm]|X[l,m]|},hlines,vlines,hspan=minimal,row{1-3}={bg=headgray},rowsep=1.5pt,colsep=3pt}
Jadwal Pelaksanaan: & Minggu 1--17 sesuai RPS. Setiap evaluasi memiliki RTM dan rubrik multi-kriteria. \\
Lain-Lain Yang Diperlukan: & LMS, modul ajar, lembar kasus, dan perangkat presentasi. \\
Pustaka: & Mengacu pustaka utama dan pendukung pada bagian RPS. \\
\end{tblr}
